\section{Precisión exponencial en el tamaño físico}
En un menor real difuso de tamaño $m\in[(1-a)j,j]$, escribimos $B=D'-\mathbf1\mathbf1^T/m$
y $d(y)=\det(I-y^2 B^TB)^{-1/2}$. Para $y$ en un intervalo fijo alrededor de uno se tiene
$F_y(w)=R_B(yw)/d(yw)$. Por R1, $|[w^i]F_y|\le e^{C\sqrt m}R^{-i}$ para un $R>1$.
Elegimos $1<r_1<R$ en los discos de R1 y de las contracciones. Para $Q=\lfloor\epsilon j\rfloor$,
$|u|\le c_0/Q$ y $1\le l\le Q$, la parte 02 de la entrega a dos tercios da
\[
 \sup_{|w|=r_1}|V_l(yw,u)|\le(Cl|u|)^l .
\]
En efecto $l=E-v$, $E\le3l$, $v\le2l$; la imposición de índices distintos cuesta a lo sumo
$\prod_{i<v}(1+i|u|)\le e^{Cv}$, y el número de contracciones dividido por $v!$
cuesta $C^E E^{E-v}$. Estas cotas valen en la bola compleja archivada.
La estimación del coeficiente del cociente físico se usa sólo con matrices reales.
\paragraph{Coincidencia exacta de coeficientes.}
Al tamaño físico $u=1/m$, cada arista aporta un propagador $H(yw)=O(w)$.
A grado $i$ sólo intervienen grafos con $E\le i$, así que la identidad formal de Wick
de R1 contiene, hasta grado $Q$, exactamente los términos de
$1+\sum_{1\le l\le Q}V_l(yw,1/m)$. Esta igualdad es coeficiente a coeficiente.
La parte retenida es acotada en $|w|=r_1$ al tomar $c_0$ pequeño y $Q/m\le c_0/2$.
La suma de las dos colas de Cauchy en $w=1$ es a lo sumo
$C(e^{C\sqrt m}+1)r_1^{-Q}$. Como $m\le j$ y $\epsilon>0$ es fijo, es $O(e^{-c_\epsilon j})$.
No se necesita una cota separada en exceso para grafos físicos con $l>Q$.
\paragraph{La integral física y el logaritmo complejo.}
Póngase $S_Q=d(y)(1+\sum_{l=1}^QV_l)$. Es holomorfa y acotada para $|y-1|<d_0$
uniformemente en la bola de datos y en $|u|\le c_0/Q$. Sea $X=m/T$, $T$ gamma de
forma $m+1$ y tasa uno. En el exterior de $I=\{|X-1|\le d_0/2\}$, R1 da
$\EE[|R_B(X)|;I^c]\le e^{-cm+C\sqrt m}$, ajustando las constantes al intervalo.
Si se reduce el intervalo anterior, en la franja retirada se usa la cota
$|R_B(y)|\le Ce^{C\sqrt m}$ y la probabilidad gamma exponencial.
Sea $\mu_i=\EE(X-1)^i$ en $u=1/m$. La recurrencia racional
\[
 \mu_0=1,\quad \mu_1=0,\qquad
 (1-iu)\mu_{i+1}=iu(2\mu_i+\mu_{i-1})
\]
se obtiene por integración por partes. Por inducción hasta $4Q+4$ se tiene
$|\mu_i(u)|\le(C\sqrt{(4Q+4)|u|})^i$. Definimos
$H_Q(u)=\sum_{i=0}^{2Q+1}\mu_i(u)\partial_y^i S_Q(1,u)/i!$.
El resto de Taylor en $I$ se integra usando el momento par $2Q+2$ y cuesta
$(C Q/m)^{Q+1}$. Para el polinomio en $I^c$, Cauchy--Schwarz da por monomio
$\Prb(I^c)^{1/2}(\mu_{2i}(1/m))^{1/2}$ antes de multiplicar por la cota de Cauchy.
Primero se reduce $c_0$ y luego $\epsilon$; ambas contribuciones son $O(e^{-cj})$.
En la bola compleja reducida $d(1)$ está uniformemente separado de cero, con rama acotada.
Por la cota de $V_l$ y las cotas de momentos, $|H_Q-d(1)|<|d(1)|/2$ en $|u|\le c_0/Q$.
La rama $\mathcal L_Q=\log H_Q$ es acotada y holomorfa en esos datos. Por la mezcla exacta,
\[
 \log(m^m\per D'/m!)=\mathcal L_Q(1/m,d')+O(e^{-c j}).
\]
La comparación con el permanente se afirma sólo en los puntos físicos.
